\documentclass[12pt,letterpaper]{article}

\usepackage[spanish]{babel}
\usepackage[utf8]{inputenc}
\usepackage[T1]{fontenc}
\usepackage{graphicx}
\usepackage{fancyhdr}
\usepackage{geometry}
\usepackage{amsmath, amssymb}
\usepackage{titlesec}


\geometry{margin=2.5cm}
\setlength{\headheight}{55pt}


\pagestyle{fancy}
\fancyhf{}

\fancyhead[L]{
    \includegraphics[height=1.5cm]{ur_image.png}
}

\fancyhead[R]{
\begin{tabular}{r}
\textbf{Monitoria Álgebra Lineal} \\
Gabriel Salcedo\\
\end{tabular}
}

\renewcommand{\headrulewidth}{0.4pt}

\begin{document}
\section*{Determinantes}

Calcule el determinante de las siguientes matrices.

\begin{enumerate}

\item Calcule el determinante de la matriz usando la regla de Sarrus.

\[
A=
\begin{pmatrix}
2 & -1 & 3\\
0 & 4 & 5\\
1 & 2 & 6
\end{pmatrix}
\]

\item Calcule el determinante de la siguiente matriz utilizando expansión de Laplace por la primera fila.

\[
C=
\begin{pmatrix}
1 & 2 & 0\\
3 & -1 & 4\\
2 & 5 & 6
\end{pmatrix}
\]


\item Calcule el determinante de la matriz.

\[
E=
\begin{pmatrix}
1 & 2 & 3 & 0\\
0 & 4 & 5 & 1\\
2 & 0 & 1 & 3\\
1 & 2 & 0 & 4
\end{pmatrix}
\]

(Sugerencia: utilice expansión de Laplace en la fila o columna que considere más conveniente).

\item Determine el valor de $k$ para el cual el determinante de la matriz es cero.

\[
F=
\begin{pmatrix}
1 & 2 & 3\\
2 & k & 4\\
1 & 0 & 2
\end{pmatrix}
\]

\end{enumerate}

\vspace{0.3cm}
\section*{Regla de Cramer}

Resuelva los siguientes sistemas de ecuaciones lineales utilizando la \textbf{Regla de Cramer}.

\begin{enumerate}


\item Resuelva el sistema de ecuaciones:

\[
\begin{cases}
x + 2y = 7 \\
3x + y = 5
\end{cases}
\]


\item Resuelva el sistema:

\[
\begin{cases}
2x + y - z = 3 \\
x + 3y + 2z = 13 \\
3x - y + z = 4
\end{cases}
\]

\item Determine el valor de $k$ para el cual el sistema tiene solución única.

\[
\begin{cases}
kx + y = 4 \\
2x + 3y = 7
\end{cases}
\]

(Sugerencia: calcule el determinante de la matriz de coeficientes).

\item Determine si el sistema puede resolverse con la Regla de Cramer.

\[
\begin{cases}
x + 2y + z = 4 \\
2x + 4y + 2z = 8 \\
x + y + z = 3
\end{cases}
\]

Justifique su respuesta calculando el determinante de la matriz de coeficientes.

\end{enumerate}
\vspace{0.3cm}
\section*{Matriz Inversa}

En los siguientes ejercicios, determine la matriz inversa si existe.

\begin{enumerate}

\item Halle la inversa de la matriz:

\[
A=
\begin{pmatrix}
2 & 1 \\
5 & 3
\end{pmatrix}
\]

\item Halle la inversa de la matriz:

\[
B=
\begin{pmatrix}
4 & -1 \\
2 & 1
\end{pmatrix}
\]

\item Determine si la matriz es invertible y encuentre su inversa si existe.

\[
D=
\begin{pmatrix}
2 & 1 & 0\\
1 & 3 & 2\\
0 & 1 & 1
\end{pmatrix}
\]

\item Determine si la matriz tiene inversa.

\[
E=
\begin{pmatrix}
1 & 2 & 3\\
2 & 4 & 6\\
1 & 0 & 1
\end{pmatrix}
\]

Justifique su respuesta calculando el determinante.

\item Utilice la matriz inversa para resolver el siguiente sistema de ecuaciones:

\[
\begin{cases}
2x + y = 5 \\
x + 3y = 7
\end{cases}
\]

(Sugerencia: escriba el sistema como $AX=B$ y calcule $X=A^{-1}B$).

\item Utilice la matriz inversa para resolver el sistema:

\[
\begin{cases}
x + y + z = 6 \\
2x + 3y + z = 10 \\
x + 2y + 3z = 14
\end{cases}
\]

\end{enumerate}

\end{document}
